\documentclass[a4paper,11pt]{article}
\usepackage[utf8]{inputenc}
\usepackage{amsmath}
\usepackage{amssymb}
\usepackage{amsthm}
\usepackage{amsfonts}
\usepackage{amsbsy}
\usepackage{esint}
\usepackage{mathtools}
\usepackage{array}
\usepackage{tikz}
\usepackage{url}
\usepackage{graphicx}
\usepackage{physics}
\usepackage{dsfont}
\usepackage{xcolor}
\usepackage[colorlinks=true, allcolors=blue]{hyperref}
\newcommand{\Schatten}{{\mathfrak S}}
\newcommand{\per}{\mathrm{per}}
\newcommand{\Ran}{\mathrm{Ran}}
\newcommand{\Trper}{\underline{\Tr}}
\newcommand{\BZ}{\mathcal{B}}
\theoremstyle{plain}
\newtheorem{theorem}{Theorem}[section]
\newtheorem{definition}[theorem]{Definition}
\newtheorem{assumption}[theorem]{Assumption}
\newtheorem{lemma}[theorem]{Lemma}
\newtheorem{corollary}[theorem]{Corollary}
\newtheorem{proposition}[theorem]{Proposition}
\newtheorem{hyp}[theorem]{Assumption}
\theoremstyle{remark}
\newtheorem{remark}[theorem]{Remark}
\newtheorem*{notation}{Notation}
\newtheorem{example}[theorem]{Example}
\renewcommand{\vec}[1]{#1}
\usepackage{tikz-cd}

\makeatletter
\tikzset{
  column sep/.code=\def\pgfmatrixcolumnsep{\pgf@matrix@xscale*(#1)},
  row sep/.code   =\def\pgfmatrixrowsep{\pgf@matrix@yscale*(#1)},
  matrix xscale/.code=
    \pgfmathsetmacro\pgf@matrix@xscale{\pgf@matrix@xscale*(#1)},
  matrix yscale/.code=
    \pgfmathsetmacro\pgf@matrix@yscale{\pgf@matrix@yscale*(#1)},
  matrix scale/.style={/tikz/matrix xscale={#1},/tikz/matrix yscale={#1}}}
\def\pgf@matrix@xscale{1}
\def\pgf@matrix@yscale{1}
\makeatother


\usepackage[margin=2.0cm]{geometry}


\def\R{{\mathbb R}}
\def\N{{\mathbb N}}
\def\C{{\mathbb C}}
\def\Z{{\mathbb Z}}
\def\eps{{\varepsilon}}

\def\cB{{\mathcal B}}
\def\cC{{\mathcal C}}
\def\cH{{\mathcal H}}
\def\cL{{\mathcal L}}
\def\cN{{\mathcal N}}
\def\cR{{\mathcal R}}


\def\1{{\mathds 1}}

\def\dps{\displaystyle }


\newcommand{\sub}[1]{_{\mathrm{#1}}}
\newcommand{\su}[1]{^{\mathrm{#1}}}


\title{Conductivity at conical crossings}
\author{}
\date{}
\begin{document}
\maketitle

\begin{abstract}
We consider independent electrons in a periodic crystal in their ground
state, and turn on a uniform electric field at some prescribed time.
We rigorously define the current per unit area and study the static
conductivity of a two-dimensional semimetal. The conductivity tensor is isotropic and each
conical intersection contributes as $\frac{1}{16}$ to the total conductivity.
\end{abstract}
\tableofcontents
\section{Introduction}

We consider a two-dimensional perfect crystal with periodic lattice
$\cR\subset\R^2$. We use atomic units and a spinless model, with Hamiltonian
\begin{equation}\label{eq:periodic_Hamiltonian}
H=-\frac12\Delta+V\quad\hbox{on }L^2(\R^2;\C),
\end{equation}
where $V\in H^1_{\rm per}(\Omega;\R)$ and $\Omega$ is a unit cell of $\cR$.
At zero temperature, the ground-state density matrix is given by
\begin{equation}\label{eq:GSDM}
\gamma_0=\mathds 1(H\le\mu_{\rm F}),
\end{equation}
where $\mu_{\rm F}\in\R$ is the Fermi level.

The purpose of this article is to analyze mathematically the behavior
of the current appearing in the crystal when a uniform external field
is turned on instantaneously at the initial time $t=0$.
Fix $e_\beta,e_\alpha\in\R^2$ once and for all, for the field and observation
directions respectively. They need not be normalized, linearly independent,
or orthogonal. Put $x_\beta=x\cdot e_\beta$ and $p=-i\nabla$. For $\eps\ge0$,
\begin{align}\label{def:Heps}
H_\eps=H-\eps x_\beta,
\end{align}
where $H_\eps$ denotes the self-adjoint closure of its restriction to
$C_c^\infty(\R^2;\C)$ (see Proposition~\ref{prop:setting}). The state at
$t\ge0$ is
\begin{equation}\label{def:gammaeps}
\gamma_\eps(t)=e^{-itH_\eps}\gamma_0e^{itH_\eps}.
\end{equation}
The current is defined as
\begin{equation}\label{eq:current}
j_\eps(t)=\Trper(J\gamma_\eps(t)),
\end{equation}
where $\Trper$ is the trace per unit area, defined in
Section~\ref{sec:notation}, and
\begin{equation}\label{def:J}
J=p\cdot e_\alpha.
\end{equation}


For time-reversible 2D semimetals such as graphene, the time-averaged conductivity
\begin{align}\label{def:conductivity}
\sigma = \lim_{t\to\infty}\frac1t\int_0^t
\lim_{\eps\to0^+}\frac{j_\eps(t')}{\eps}\,dt'
\end{align}
has a finite value equal to $\frac1{16}e_\alpha\cdot e_\beta$ times
the number of Dirac points in the Brillouin zone.


Note that our results use an averaging in time, and we are unable to
conclude anything about what would be the naive definition of the
conductivity
\begin{align}
  \label{eq:conductivity_naive}
  \lim_{t \to  \infty} \lim_{\eps \to 0^+} \frac{j_\eps(t)}{\eps}.
\end{align}
A form of averaging of time fluctuations is always necessary to infer
zero-frequency behavior from step responses in non-dissipative
systems, even in the linear case. The easiest way to see this is by
the very simple model for the relationship between an input $I(t)$ and
an output $O(t)$:
  \begin{equation} \label{eq:toymodel}
    i \dot O(t) = \omega O(t) + I(t).
  \end{equation}
  This simplified model describes a forced oscillator with eigenfrequency $\omega$, and arises from the linear response of the time-dependent Schr\"odinger equation of a two-level system. For a
  constant input~$I_{0}$, there is a steady state solution
  $O_0 = \widehat R_0 I_{0}$, where $\widehat R_0 = -\frac 1 \omega$ is
  the zero-frequency transfer function of the system. However, since
  this system is oscillatory, this steady state may never be reached:
  if $I$ is brutally switched on at time $0$ with $I(t) = O(t) = 0,$
  if $t \leq 0$, $I(t) = I_{0}$ if $t > 0$, then
  $ O(t) = O_0(1- {\rm e}^{-i\omega t}) = \widehat{R}_0 I_{0} (1- {\rm
    e}^{-i\omega t})$ and we cannot define $\widehat R_0$ as the limit
  of $O(t)/I_{0}$ when $t$ goes to infinity. However, by averaging we
  obtain
  \begin{align*}
     \lim_{t\to\infty}\frac 1 t  \int_{0}^{t} \frac{O(t')}{I_{0}} \, dt' = \widehat R_0.
   \end{align*}

In our results in the linear response regime, we consider the quantity
\begin{align*}
  j^{\rm LR}(t) = \lim_{\eps \to 0^+} \frac {j_\eps(t)}{\eps}.
\end{align*}
and then compute the conductivity in the limit $t \to \infty$. This
order of limits is necessary to ensure that the electronic state never
deviates significantly from its equilibrium. Considering the opposite
limit, i.e. the infinite-time dynamics of
$j_\eps(t)$ at $\eps$ fixed, is an extremely
hard problem, as it is for any dynamics of non-dissipative systems.
Our method of proof is based on the standard gauge change
$\widetilde\psi(x,t)=e^{-i\eps t x_\beta}\psi(x,t)$
that transforms the constant in time but non-spatially-periodic
Hamiltonian $H_\eps$ into the time-dependent Hamiltonian
$\widetilde H_\eps(t)=\frac12(p+\eps t e_\beta)^2+V$.
This Hamiltonian is spatially periodic, and the study of its dynamics
can be reduced via Bloch--Floquet theory to that of its fibers
$H_{k+\eps t e_\beta}$ acting on periodic functions
(Section~\ref{sec:preliminaries}). Fiber by fiber, this time-dependent
Hamiltonian can then be treated using the tools of time-dependent
perturbation theory (Section~\ref{sec:linear-response}). For each fixed $t$,
we use linear response to expand $j_\eps(t)$ to first order in $\eps$
(Proposition~\ref{prop:Dyson}).


The structure of the paper is as follows. We describe our results in
Section~\ref{sec:results}: we define the current in
Proposition~\ref{prop:setting}, and state Theorem~\ref{theo:semi-metal}.
We devote Section~\ref{sec:preliminaries} to preliminaries about the
regularity and Bloch decomposition of the current.
Section~\ref{sec:linear-response} states and proves the linear response result.
Section~\ref{sec:semimet} is devoted to the proof of the conductivity
formula. Finally two short Appendices are devoted to technical issues.

\section{Main results: electrical current in periodic materials}\label{sec:results}


\subsection{Notation}
\label{sec:notation}
We denote by $\cR^*$ the dual lattice
\[
\cR^*=\{K\in\R^2:K\cdot R\in2\pi\Z\hbox{ for all }R\in\cR\},
\]
by $\Omega$ a unit cell in physical space, and by $\cB$ a unit cell in
reciprocal space, called the Brillouin zone. We write
$\partial_\alpha=e_\alpha\cdot\nabla_k$ and
$\partial_\beta=e_\beta\cdot\nabla_k$ for the two fixed directional derivatives.



The {\it $\cR$-periodic Lebesgue and
Sobolev spaces} are denoted by
\begin{align*}
L^p_{\rm per}&:=\left\{ u \in L^p_{\rm loc}(\R^2;\C) \; | \; u \mbox{ $\cR$-periodic} \right\}, \\
H_{\rm per}^{s}&:=\left\{ u \in H_{\rm loc}^{s}(\R^2;\C) \; | \; u \mbox{ $\cR$-periodic} \right\}.
\end{align*}
The space of bounded linear operators on a Hilbert space~$\mathcal{H}$ is denoted by ${\mathcal L}(\mathcal{H})$, and the {Schatten class} of bounded operators~$A \in {\mathcal L}(\mathcal{H})$ such that $\Tr(|A|^p) < \infty$ by $\Schatten_{p}(\mathcal{H})$.
For $R \in \cR$, we denote by~$\tau_R$ the {translation operator} formally defined by $\tau_R \phi = \phi(\cdot - R)$. Depending on the context, $\tau_R$ will be seen as a unitary operator on $L^2(\R^2;\C)$, or as a linear operator on some $\cR$-translation invariant subspace of ${\mathcal D}'(\R^2;\C)$.
A bounded operator on $L^2(\R^2;\C)$ is called $\cR$-periodic if it commutes with~$\tau_R$ for all $R \in \cR$. An unbounded self-adjoint operator on $L^2(\R^2;\C)$ is called $\cR$-periodic if its resolvent is $\cR$-periodic. A bounded $\cR$-periodic operator $A \in {\mathcal L}(L^2(\R^2;\C))$ is called locally trace-class if $\chi A \chi \in \Schatten_{1}(L^2(\R^2;\C))$ for any compactly supported function $\chi \in L^\infty(\R^2;\C)$. For $p \ge 1$, we denote by $\Schatten_{p,\per}$ the space of $\cR$-periodic operators $A \in {\mathcal L}(L^2(\R^2;\C))$ such that $|A|^p$ is locally trace class. Any operator $A \in \Schatten_{1,\per}$ has a density $\rho_A \in L^1_{\rm per}$ characterized by
$$
\forall \chi \in C^\infty_{\rm c}(\R^2;\C), \quad \Tr(A\chi) = \int_{\R^2} \rho_A \chi.
$$
The trace per unit area of an operator $A \in \Schatten_{1,\per}$ is defined as
$$
\Trper(A) = \frac{1}{|\Omega|} \Tr_{L^2(\R^2;\C)}(\1_\Omega A \1_\Omega) =
\fint_{\Omega}  \rho_A,
$$
where $\1_\Omega$ is the characteristic function of the unit cell
$\Omega$, and $\fint_\Omega$ is a shorthand notation for $\frac 1
{|\Omega|} \int_{\Omega}$. This formula is independent of the choice
of the unit cell $\Omega$.


\medskip

Since we are dealing here with periodic materials, we will use the {Bloch transform} (also
called Bloch-Floquet transform) \cite{ReedSimon4}. For $K \in \cR^\ast$, let $T_K$ be the unitary multiplication operator on $L^2_{\rm per}$ defined by
$$
\forall v \in L^2_{\rm per}, \quad (T_Kv)(x) = {\rm e}^{-iK\cdot x} v(x) \quad \mbox{ for a.a. } x \in \R^2,
$$
and
\begin{align*}
&L^2_{\rm qp}(L^2_{\rm per})\\
&:=\left\{ \R^2 \ni k \mapsto u_k \in L^2_{\rm per} \; | \; \int_\cB \|u_k\|_{L^2_{\rm per}}^2 \, dk < \infty, \; u_{k+K}=T_K u_k \mbox{ for all } K \in \cR^\ast \mbox{ and a.a. } k \in \R^2 \right\},
\end{align*}
the Hilbert space of $\cR^\ast$-quasi-periodic $L^2_{\rm per}$-valued functions on $\R^2$ endowed with the inner product
$$
\langle u,v\rangle_{L^2_{\rm qp}(L^2_{\rm per})}= \fint_\cB \langle u_k,v_k\rangle_{L^2_{\rm per}} \, dk.
$$
Here and below, the subscript qp refers to the quasi-periodicity property.
The Bloch transform then is the unitary map from $L^2(\R^2;\C)$ to $L^2_{\rm qp}(L^2_{\rm per})$ defined for $u \in C^\infty_{\rm c}(\R^2;\C)$ by
\begin{equation}\label{eq:Bloch_transform}
\forall k \in \R^2, \quad \forall x \in \R^2, \quad u_k(x) = \sum_{R \in \cR} u(x+R) \, {\rm e}^{-ik \cdot (x+R)}.
\end{equation}
Its inverse is given by
\begin{equation}\label{eq:inverse_Bloch_transform}
u(x) =\fint_\cB {\rm e}^{ik\cdot x} u_{k}(x) \, dk, \quad \mbox{for a.a. } x \in \R^2.
\end{equation}
Any $\cR$-periodic
operator $A \in {\mathcal L}(L^2(\R^2;\C))$ is decomposed by the Bloch
transform in the sense that there exists a function
$k \mapsto A_k$ in $L^\infty_{\rm qp}({\mathcal L}(L^2_{\rm per}))$ such that for any
$u \in L^2(\R^2;\C)$ and almost all $k \in \R^2$, $(Au)_k = A_ku_k$, and
\begin{equation}\label{eq:qpBloch}
A_{k+K} = T_K A_k T_K^\ast, \quad \mbox{ for all } K \in \cR^\ast \mbox{ and a.a. } k \in \R^2.
\end{equation}
The $A_{k}$'s are called the fibers of the operator $A$.
If $A \in \Schatten_{1,\per}$, then the function $k \mapsto A_{k}$ is
in $L^1_{\rm qp}(\Schatten_{1}(L^2_{\rm per}))$, the function $k
\mapsto \Tr(A_k)$ is in $L^1_{\rm loc}(\R^2)$, $\cR^{*}$-periodic, and we have
$$
\Trper(A) =(2\pi)^{-2} \int_\cB \Tr(A_k) \, dk.
$$
The Bloch decomposition theorem can be extended to unbounded
$\cR$-periodic self-adjoint operators using the resolvent \cite{ReedSimon4}.

\medskip


In the case of the periodic Hamiltonian operator $H$ given by \eqref{eq:periodic_Hamiltonian}, we have
  \begin{align} \label{eq:Hk}
    H_k = \frac 1 2 (-i\nabla+k)^2 + V.
  \end{align}
  For each $k \in \R^2$, $H_k$ is a bounded below self-adjoint operator on $L^{2}_{\per}$ with domain $H_{\per}^{2}$ and compact resolvent.
  Let $(\lambda_{n,k})_{n \in \N^\ast}$ be the non-decreasing sequence of eigenvalues of $H_k$ counting multiplicities
  $$
  \lambda_{1,k}\leq \lambda_{2,k}\leq \lambda_{3,k} \leq \cdots , \quad \lim_{n \to \infty} \lambda_{n,k}=+\infty,
  $$
  and we use the convention $\lambda_{0,k} = -\infty$. We denote by
  $(u_{n,k})_{n \in \N^\ast} \in (H_{\per}^{2})^{\mathbb N^{*}}$ an $L^2_{\rm per}$-orthonormal basis of associated eigenfunctions:
  \begin{align*}
    H_k u_{n,k} = \lambda_{n,k} u_{n,k}, \quad \langle u_{m,k},u_{n,k}\rangle_{L^2_{\rm per}}=\delta_{m,n}.
  \end{align*}
For $n\ge1$ and $k\in\R^2$, set
\[
P_{n,k}=\mathds 1(H_k\le\lambda_{n,k}),\qquad P_{0,k}=0.
\]
Whenever $\lambda_{n,k}<\lambda_{n+1,k}$, $P_{n,k}$ is the spectral
projector associated with the lowest $n$ eigenvalues, counting multiplicities:
\begin{equation}\label{eq:def_PNk}
P_{n,k}=\sum_{m=1}^n|u_{m,k}\rangle\langle u_{m,k}|.
\end{equation}

  Since $H_k$ is quasi-periodic, so is $P_{n,k}$, and the eigenvalues $\lambda_{n,k}$ are
  ${\mathcal R}^\ast$-periodic functions of $k$. By a min-max argument
  (see e.g. \cite{ReedSimon4}), there exists
  $\underline{C_1}, \overline{C_1} \in \R$, and
  $\underline{C_2},\overline{C_2}>0$ such that
  \begin{align}
    \label{eq:asymptotics_eigenvalues}
 \underline{C_1}+ \underline{C_2} n  \leq \lambda_{n,k} \leq  \overline{C_1}+ \overline{C_2} n.
\end{align}
The ground-state density matrix $\gamma_0$ has Bloch fibers
\begin{equation}\label{eq:gk0}
P_k=\mathds 1(H_k\le\mu_{\rm F}).
\end{equation}
Their ranks are uniformly bounded by the eigenvalue estimate above.
The current operator $J$ is also $\cR$-periodic, with fibers
\[
J_k=(p+k)\cdot e_\alpha=\partial_\alpha H_k.
\]


Lastly, for each $q \in \R^2$, we denote by $G_q$ the unitary multiplication operator on $L^2(\R^2;\C)$ defined by
\begin{equation}\label{eq:defGq}
\forall u \in L^2(\R^2;\C), \quad (G_qu)(x) = {\rm e}^{iq\cdot x} u(x) \quad \mbox{ for a.a. } x \in \R^2.
\end{equation}
The operator $G_q$ is not $\cR$-periodic, except
when $q \in \cR^\ast$ (in which case~$G_{q}$ is fibered, with
$G_{q,k} = T_{-q}$ for all $k$). However, for any $\cR$-periodic
operator $A \in {\mathcal L}(L^2(\R^2;\C))$ and any $q \in \R^2$,
the operator $G_qAG_q^\ast$ is $\cR$-periodic and its Bloch
decomposition is given by
\begin{equation}\label{eq:fibers_eiqx}
(G_qAG_q^\ast)_k = A_{k-q}, \quad \mbox{ for a.a. } k \in \R^2.
\end{equation}


\subsection{The Bloch theorem}
Before attacking the well-posedness of the current
$j_\eps(t) = \underline {\rm Tr}(J
\gamma_\eps(t))$ for $\eps, t \neq 0$, we first study an easier special case.
\begin{proposition}[Bloch theorem]\label{prop:null_current} The current satisfies
\begin{align*}
& j_0(t) = 0, \quad \forall t \ge 0 \;\;\; \mbox{(no current in the absence of external field)}, \\
& j_\eps(0)=0, \quad \forall \eps \ge 0 \;\;\;  \mbox{(continuity of the current at $t=0$).}
\end{align*}
\end{proposition}
\begin{proof}
We have
\[
j_0(t)=j_\eps(0)=\frac1{4\pi^2}\int_\mathcal B
\Tr(\partial_\alpha H_kP_k)\,dk.
\]
Since $V$ is real, $H_{-k}=\overline{H_k}$ and $P_{-k}=\overline{P_k}$,
where the bar denotes complex conjugation. Moreover,
$\partial_\alpha H_{-k}=-\overline{\partial_\alpha H_k}$.
Thus the real-valued, $\cR^*$-periodic function
$k\mapsto\Tr(\partial_\alpha H_kP_k)$ is odd, and its integral vanishes.
\end{proof}


\subsection{Definition of the current}
For $\eps>0$, the operator $H_\eps=H-\eps x_\beta$ is not
$\cR$-periodic, and we would naively expect that the density matrix
$\gamma_\eps(t)$ is not either. Yet, this operator is in fact
$\cR$-periodic. Physically, this is due to the fact that although
$-\eps x\cdot e_\beta$ is not periodic, its negative gradient is constant,
hence periodic. The proof relies on the standard gauge transform
\begin{align}\label{eq:gauge_change}
\widetilde\psi(x,t)=G_{-\eps t e_\beta}\psi(x,t)
=e^{-i\eps t x_\beta}\psi(x,t),
\end{align}
where $G_q$ was defined in~\eqref{eq:defGq}, and the gauge-transformed operators
\begin{align}\label{gauge_u}
U_\eps(t,t')=G_{-\eps t e_\beta}
 e^{-i(t-t')H_\eps}G_{-\eps t'e_\beta}^*,
\end{align}
and
\begin{equation}\label{gauge}
\widetilde\gamma_\eps(t)=G_{-\eps t e_\beta}\gamma_\eps(t)G_{-\eps t e_\beta}^*
=U_\eps(t)\gamma_0U_\eps(t)^*,
\end{equation}
where $U_\eps(t)=U_\eps(t,0)$. Through this change of gauge, the dynamics
induced by $H_\eps$ is equivalent to that induced by the time-dependent
periodic Hamiltonian
\begin{align}\label{widetilde_H}
\widetilde H_\eps(t)=G_{-\eps t e_\beta}H_\eps G_{-\eps t e_\beta}^*
+\eps x_\beta=\frac12(p+\eps t e_\beta)^2+V.
\end{align}
Its Bloch fibers are
\begin{align}\label{widetilde_H_k}
\widetilde H_{\eps,k}(t)=\frac12(p+k+\eps t e_\beta)^2+V=H_{k+\eps t e_\beta}.
\end{align}
We sum up these arguments in the proposition below, together with elements
that we shall use for defining the current.

\begin{proposition}\label{prop:setting}
Let $V\in H^1_{\rm per}(\Omega;\R)$.
\begin{enumerate}
\item For every $\eps\in\R$, $H_\eps$ is essentially self-adjoint on
$C_c^\infty(\R^2;\C)$ and therefore generates a unitary group on $L^2(\R^2;\C)$.
\item The operator $\widetilde H_\eps(t)$ is self-adjoint with domain
$H^2(\R^2;\C)$ and $\cR$-periodic. The propagator $U_\eps(t,t')$ is
$\cR$-periodic, with fibers $U_{\eps,k}(t,t')$ solving
\begin{equation}\label{eq:propagator2}
i\partial_tU_{\eps,k}(t,t')=H_{k+\eps te_\beta}U_{\eps,k}(t,t'),
\qquad U_{\eps,k}(t',t')=\mathrm{Id}.
\end{equation}
\item For $t\ge0$, $J\gamma_\eps(t)\in\Schatten_{1,\rm per}$, and
\begin{align}\label{eq:current-k}
j_\eps(t)&=\frac1{4\pi^2}\int_\cB
\Tr\left(\partial_\alpha H_{k+\eps te_\beta}\widetilde\gamma_{\eps,k}(t)\right)\,dk\\
&=\frac1{4\pi^2}\int_\cB\Tr\left(\partial_\alpha H_{k+\eps te_\beta}
U_{\eps,k}(t)P_kU_{\eps,k}(t)^*\right)\,dk.\nonumber
\end{align}
\end{enumerate}
\end{proposition}
The results of Proposition~\ref{prop:setting} are proved in
Section~\ref{sec:preliminaries} for completeness. The situation can be
summed up in the commutative diagrams of Figure~\ref{fig:diagram}.
\begin{figure}[h!]
\centering
\begin{tikzcd}[column sep=large,row sep=large]
\gamma_0\arrow{r}{e^{-itH_\eps}}\arrow{d}{\mathrm{Id}}&
\gamma_\eps(t)\arrow{d}{G_{-\eps te_\beta}}\\
\gamma_0\arrow{r}{U_\eps(t)}&\widetilde\gamma_\eps(t)
\end{tikzcd}
\medskip

\begin{tikzcd}[column sep=large,row sep=large]
\{P_k\}_k\arrow{r}{}\arrow{d}{\mathrm{Id}}&
\{\gamma_{\eps,k}(t)\}_k\arrow{d}{k\mapsto k+\eps te_\beta}\\
\{P_k\}_k\arrow{r}{\{U_{\eps,k}(t)\}_k}&
\{\widetilde\gamma_{\eps,k}(t)\}_k
\end{tikzcd}
\caption{Relations between the density matrices and their Bloch fibers.
An arrow labelled by a unitary $U$ denotes conjugation by $U$.
The right lower arrow means
$\widetilde\gamma_{\eps,k}(t)=\gamma_{\eps,k+\eps te_\beta}(t)$.}
\label{fig:diagram}
\end{figure}


\subsection{Semimetals}\label{sec:semimetal-assumptions}
The position of the Fermi level in the band diagram is key to determining
the electronic properties of the medium. We assume the following.
\begin{assumption}[semimetal]\label{ass:semimetal}
There are $N,M\ge1$ and distinct points $k_1,\ldots,k_M$ modulo $\cR^*$ such that
\[
\lambda_{N,k}<\mu_{\rm F}<\lambda_{N+1,k}\qquad
(k\notin\{k_1,\ldots,k_M\}\bmod\cR^*).
\]
For each $i=1,\ldots,M$, the two bands meet at an isolated conical crossing:
\begin{align}
\lambda_{N-1,k_i}<\lambda_{N,k_i}=\mu_{\rm F}
=\lambda_{N+1,k_i}<\lambda_{N+2,k_i},\label{hyp:semimetals_1}\\
\lambda_{N,k}=\mu_{\rm F}-v_i|k-k_i|+O(|k-k_i|^2),\label{hyp:semimetals_2}\\
\lambda_{N+1,k}=\mu_{\rm F}+v_i|k-k_i|+O(|k-k_i|^2),\label{hyp:semimetals_3}
\end{align}
where $v_i>0$.
\end{assumption}
The Fermi surface is the finite set
\begin{equation}\label{eq:Fermi_surface}
\mathcal S=\{k_1,\ldots,k_M\}\bmod\cR^*.
\end{equation}
For $k\notin\mathcal S$, $P_k=P_{N,k}$ has rank $N$. Its values at the
crossings do not affect any Brillouin-zone integral. Since $V$ is real,
$H_{-k}=\overline{H_k}$: the system has time-reversal symmetry.
The regularity assumption $V\in H^1_{\rm per}$ will be used to justify the
Dyson expansion in Proposition~\ref{prop:Dyson}.


\subsection{Conductivity}
In the following, $O(f(\eps,t))$ denotes a quantity bounded by $Cf(\eps,t)$,
where $C$ may depend on the potential, the Fermi level and the fixed
directions, but not on $t$ or $\eps$. A subscript $T$ allows the constant
to depend on a fixed time horizon $T$.
\begin{theorem}\label{theo:semi-metal}
Under Assumption~\ref{ass:semimetal}, the iterated limit exists and
\[
\sigma=\lim_{T\to\infty}\frac1T\int_0^T
\left(\lim_{\eps\to0^+}\frac{j_\eps(t)}{\eps}\right)\,dt
=\frac{M}{16}\,e_\alpha\cdot e_\beta.
\]
\end{theorem}
The value of the conductivity is universal: it does not depend on the
characteristics of the Hamiltonian, but only on the number of conical
crossings. The conductivity tensor is isotropic and each conical
intersection contributes $1/16$.


\section{Bloch decomposition and regularity of the current}\label{sec:preliminaries}
In this section, we prove Proposition~\ref{prop:setting}. We first
point out an alternative interpretation of the gauge change.
Formally, $\gamma_\eps(t)$ satisfies
\[
i\partial_t\gamma_\eps=[H,\gamma_\eps]-\eps[x_\beta,\gamma_\eps].
\]
The commutator $[x_\beta,\gamma_\eps]$ is $\cR$-periodic, with fibers
$i\partial_\beta\gamma_{\eps,k}$. Therefore the fibers satisfy
\[
i\partial_t\gamma_{\eps,k}+i\eps\partial_\beta\gamma_{\eps,k}
=[H_k,\gamma_{\eps,k}]=L_{H_k}\gamma_{\eps,k},
\]
where $L_{H_k}=[H_k,\cdot]$ is the Liouvillian associated with $H_k$
(see Section~\ref{sec:louvillian_intro}). The left-hand side is a linear
advection equation, which suggests the method of characteristics: setting
\begin{equation}\label{eq:change_of_variable}
\widetilde\gamma_{\eps,k}(t)=\gamma_{\eps,k+\eps te_\beta}(t),
\end{equation}
we obtain
\[
i\partial_t\widetilde\gamma_{\eps,k}(t)
=[H_{k+\eps te_\beta},\widetilde\gamma_{\eps,k}(t)].
\]
The gauge transform $G_{-\eps te_\beta}$, equivalent to the change of
variable~\eqref{eq:change_of_variable}, makes these remarks rigorous.
We now prove Proposition~\ref{prop:setting}.


{\bf Proof of the first assertion.}
The essential self-adjointness of $H_\eps$ follows from an extension of the Faris-Lavine theorem~\cite[Theorem~X.38]{ReedSimon2}.  Let $\mathcal{C} = C^\infty_{\rm c}(\R^2;\C)$ be the set of infinitely differentiable, compactly supported functions.
\begin{lemma}[Faris-Lavine criterion]
\label{FarisLavine}
Let $V$ and $W$ be real-valued measurable functions on $\R^2$,
with $W\in L^2_{\rm loc}(\R^2;\R)$. Suppose that
\begin{enumerate}
\item there exist $c>0$ and $f\ge0$ such that
$W(x) \geq - c|x|^2 - f$, for a.a. $x\in \R^2$;
\item
$-\frac12\Delta + V + W +2c|x|^2$ is essentially self-adjoint on $\mathcal{C}$;
\item for some $a<1$, $-\frac a2\Delta + V$ is bounded below on $\mathcal{C}$.
\end{enumerate}
Then $-\frac12\Delta + V + W$ is essentially self-adjoint on $\mathcal{C}$.
\end{lemma}
The proof of the above lemma is postponed until Appendix~\ref{sec:FL_lemma}.
\medskip


We apply Lemma~\ref{FarisLavine} with $W=-\eps x_\beta$ and $c=1$,
noting that
\[
-\eps x\cdot e_\beta\ge-|x|^2-\frac{\eps^2|e_\beta|^2}{4}.
\]
In two dimensions, $V\in L^2_{\rm per}$ is infinitesimally bounded with
respect to $-\Delta$, hence also with respect to the harmonic oscillator
$-\frac12\Delta+2|x|^2$. The linear potential $-\eps x_\beta$ is likewise
infinitesimally bounded with respect to this oscillator. The Kato--Rellich
theorem therefore gives the essential self-adjointness of
$-\frac12\Delta+V-\eps x_\beta+2|x|^2$ on the core $\mathcal C$.
Also, $-\frac a2\Delta+V$ is bounded below for every $a>0$.
Lemma~\ref{FarisLavine} gives that $H_\eps$ is essentially self-adjoint on
$\mathcal C$, and its propagator is the unitary group
$(e^{-itH_\eps})_{t\in\R}$.

\medskip

{\bf Proof of the second assertion.}
The self-adjointness of
\[
\widetilde H_\eps(t)=\frac12(p+\eps te_\beta)^2+V
\]
on $H^2(\R^2;\C)$ follows from the Kato--Rellich theorem.
To show the well-posedness of the dynamics, it suffices to study its
fibers $H_{k+\eps te_\beta}$. We use the following lemma on time-dependent
perturbations of the free-particle Hamiltonian on $L^2_{\rm per}$.


\begin{lemma}
\label{lemmaUnitaryPropagator} Let $H_{0}:=-\frac 12\Delta$ be the free-particle Hamiltonian on $L^2_{\rm per}$, and a $C^1$ map $H_{1}:[0,T]\to\mathcal L(H_{\rm per}^{2},L^2_{\rm per})$,
taking its values in the set of $H_{0}$-bounded self-adjoint operators on $L^2_{\rm per}$ with relative bound lower than $1$, that is: there exist $0<a<1$ and $b>0$ such that
\begin{align}
  \label{eq:hyp_lap_bounded}
\forall t \in [0,T],\quad \forall \phi \in H_{\rm per}^{2}, \quad \|H_{1}(t)\phi\|_{L^2_{\rm per}} \leq a \|H_{0} \phi \|_{L^2_{\rm per}} + b \|\phi\|_{L^2_{\rm per}}.
\end{align}
Then, for all $t \in [0,T]$, the operator defined by
 $H(t) = H_{0} + H_{1}(t)$
 is self-adjoint on $L^2_{\rm per}$ with domain $H_{\rm per}^{2}$, and there exists a unique unitary propagator $(\mathcal{U}(t))_{t\in[0,T]}$ on $L^2_{\rm per}$ such that for $t \in [0,T]$, and $\phi_0 \in H_{\rm per}^{2}$, $\phi: t \mapsto \mathcal{U}(t)\phi_0$ is in $C([0,T];H_{\rm per}^{2})\cap C^1([0,T];L^2_{\rm per})$, and solves the time-dependent Schr\"odinger equation
\[i \partial_t \phi(t) = H(t) \phi (t), \qquad \phi(0) = \phi_0.\]
\end{lemma}
The proof of the above lemma is postponed to Appendix~\ref{sec:unitaryprop}.

\medskip

For $t\in[0,T]$, write $H_{k+\eps te_\beta}=H_0+H_1(t)$, with
\[
H_1(t)=(k+\eps te_\beta)\cdot p+\frac12|k+\eps te_\beta|^2+V.
\]
The Sobolev embedding $H^2_{\rm per}\subset L^\infty_{\rm per}$ and
$V\in L^2_{\rm per}$ imply the conditions of
Lemma~\ref{lemmaUnitaryPropagator}; the result follows.


\noindent
{\bf Proof of the third assertion.} We first compute the fibers of
$\gamma_\eps(t)$. Using~\eqref{eq:fibers_eiqx}, we have
\begin{align}
\gamma_{\eps,k}(t)
&=\left(G_{-\eps te_\beta}^*U_\eps(t)\gamma_0U_\eps(t)^*G_{-\eps te_\beta}\right)_k\\
&=U_{\eps,k-\eps te_\beta}(t)P_{k-\eps te_\beta}
 U_{\eps,k-\eps te_\beta}(t)^*\label{eq:gebkt}\\
&=\sum_{n:\,\lambda_{n,k-\eps te_\beta}\le\mu_{\rm F}}
 |U_{\eps,k-\eps te_\beta}(t)u_{n,k-\eps te_\beta}\rangle
 \langle U_{\eps,k-\eps te_\beta}(t)u_{n,k-\eps te_\beta}|,\nonumber
\end{align}
for almost every $k$. Since $u_{n,k}\in H^2_{\rm per}$, we deduce that
\[
(J\gamma_\eps(t))_k=\partial_\alpha H_k\gamma_{\eps,k}(t)
\]
is trace-class and finite-rank, with a uniform bound for $k\in\cB$.
The current is therefore well-defined. Since
$k\mapsto\Tr(\partial_\alpha H_k\gamma_{\eps,k}(t))$ is $\cR^*$-periodic,
\begin{align*}
j_\eps(t)
&=\frac1{4\pi^2}\int_\cB\Tr(\partial_\alpha H_k\gamma_{\eps,k}(t))\,dk\\
&=\frac1{4\pi^2}\int_\cB\Tr(\partial_\alpha H_{k+\eps te_\beta}
\gamma_{\eps,k+\eps te_\beta}(t))\,dk,
\end{align*}
which gives~\eqref{eq:current-k}.


\section{Linear response}\label{sec:linear-response}

\subsection{The Liouvillian and its partial inverse}
\label{sec:louvillian_intro}
In order to formulate and interpret our results, it is convenient to make use of the Liouvillian formalism. For each $k\notin\mathcal S$, the Liouvillian
$L_k$ associated with $H_k$ is the commutator
\begin{equation}\label{eq:Liouvillian}
L_k A=[H_k,A].
\end{equation}
The operator $L_k$ is to density matrices what the Hamiltonian $H_k$ is
to pure states: its evolution is
\begin{equation}\label{eq:Liouvillian2}
e^{-itL_k}A=e^{-itH_k}Ae^{itH_k}.
\end{equation}
More precisely, $L_k$ is the self-adjoint generator of this
strongly-continuous unitary group on the Hilbert--Schmidt operators.
Its action is easily understood in the eigenbasis of $H_k$:
\[
L_k|u_{n,k}\rangle\langle u_{m,k}|=
(\lambda_{n,k}-\lambda_{m,k})|u_{n,k}\rangle\langle u_{m,k}|.
\]
The operator $L_k$ is not invertible. However, it is invertible
on the block off-diagonal operators, i.e. operators $A$ such that
$P_kAP_k=(1-P_k)A(1-P_k)=0$. Its partial inverse $L_k^+$ is given by
\begin{align}\label{eq:explicit_inverse_liouvillian}
L_k^+A=\sum_{n\le N<m}
\frac{\langle u_{m,k},Au_{n,k}\rangle
 |u_{m,k}\rangle\langle u_{n,k}|
-\langle u_{n,k},Au_{m,k}\rangle
 |u_{n,k}\rangle\langle u_{m,k}|}
{\lambda_{m,k}-\lambda_{n,k}}.
\end{align}
The gap and the finite number of occupied states ensure convergence in
Hilbert--Schmidt norm for bounded $A$. We extend $L_k^+$ to all bounded
operators by setting it to zero on the diagonal blocks. In particular,
$L_kL_k^+A=A$ for off-diagonal $A$.

\subsection{Spectral derivatives and propagation bounds}
For $k\notin\mathcal S$, the contour formula
\begin{equation}\label{eq:int_rep}
P_k=\frac{1}{2\pi i}\oint_{\mathcal C}(z-H_k)^{-1}\,dz,
\end{equation}
where $\mathcal C$ encloses the occupied eigenvalues, gives a locally
analytic projector. Differentiating $[H_k,P_k]=0$, we get
\[
L_k\partial_\beta P_k=[P_k,\partial_\beta H_k].
\]
Both sides are off-diagonal (for $\partial_\beta P_k$, differentiate
$P_k^2=P_k$). Hence
\begin{equation}\label{eq:Pprime}
\partial_\beta P_k=L_k^+[P_k,\partial_\beta H_k].
\end{equation}

We also need propagation bounds that do not require a spectral gap.
Put $h_\eps(t)=H_{k+\eps te_\beta}+\mu$, with
$\mu=1-\min\sigma(H)$, so that $h_\eps(t)\ge1$.
For $0\le t\le T$, $0\le\eps\le1$ and $k\in\mathcal B$,
the graph and form norms of $h_\eps(t)$ are uniformly equivalent to
the $H^2_{\rm per}$ and $H^1_{\rm per}$ norms.
Let $\psi(t)=U_{\eps,k}(t,t')\psi$ for $\psi\in H^2_{\rm per}$.
Differentiating the energy, we get
\[
\frac{d}{dt}\langle\psi(t),h_\eps(t)\psi(t)\rangle
=\eps\langle\psi(t),\partial_\beta H_{k+\eps te_\beta}\psi(t)\rangle.
\]
The absolute value of the right-hand side is bounded by
$C_T\langle\psi(t),h_\eps(t)\psi(t)\rangle$.
By the Gr\"onwall lemma and the equivalence of form norms, we obtain
\begin{equation}\label{eq:fixed_time_propagation}
\left\|(H_k+\mu)^{1/2}U_{\eps,k}(t,t')
(H_k+\mu)^{-1/2}\right\|\le C_T,
\qquad 0\le t'\le t\le T,\quad 0\le\eps\le1.
\end{equation}
The estimate extends by density to initial data in $H^1_{\rm per}$.
By quasi-periodicity, it is uniform for all $k\in\R^2$.

\subsection{Linear response}
\label{subsec:lr}

We now aim at obtaining an expansion of the current to first order in
$\eps$ for a given $t$, based on a Dyson expansion.
This is a classical computation in response theory, sometimes known
as the Kubo formula. This gives
a remainder that does not depend on a gap, and will therefore be
useful for the study of semimetals near Dirac points.
\begin{proposition}\label{prop:Dyson}  Let $H$ be the periodic Hamiltonian defined in~\eqref{eq:periodic_Hamiltonian}. For every $T>0$, there exists a constant $C_T$ such that for all $k\notin\mathcal S$, $0\le t\le T$ and $0\le\eps\le1$,
\begin{align}
 {\rm Tr} \left(\partial_\alpha H_{k+\eps e_\beta t}U_{\eps,k}(t)\, P_k\, U_{\eps,k}(t)^*\right) =& \partial_\alpha {\rm Tr} \left(H_k P_k\right) + \eps t  {\partial_\alpha\partial_\beta} \left( \Tr (H_k P_k) \right) \nonumber  \\ & - i \eps \Tr(\partial_\alpha H_k (e^{-i t L_k} - 1) L_{k}^{+}  \partial_\beta P_k) + \rho^\eps_k(t),
\end{align}
with
\begin{equation}\label{eq:boundrho}
|\rho^\eps_k(t)| \le C_T \eps^2.
\end{equation}
\end{proposition}

\begin{proof}[Proof]  Let $k$ be such that $\lambda_{N+1,k}-\lambda_{N,k} > 0$. Since $k' \mapsto {\rm Tr} \left(H_{k'} P_{k'}\right)$ is real-analytic in a neighborhood of $k$, we have by Hellmann-Feynman theorem
\begin{align*}
 \partial_\alpha {\rm Tr} \left(H_k P_k\right) =  {\rm Tr} \left(\partial_\alpha H_k P_k\right) \quad \mbox{and} \quad
  {\partial_\alpha\partial_\beta} {\rm Tr} \left(H_k P_k\right) =  {\rm Tr}
\left(  {\partial_\alpha\partial_\beta}
 H_k P_k\right) +  {\rm Tr} \left(\partial_\alpha H_k \partial_\beta P_k\right).
\end{align*}
We also have $\partial_\alpha H_{k+\eps e_\beta t} =
\partial_\alpha H_k + \eps t e_\alpha \cdot e_\beta$. It follows that
\begin{align*}
\rho^\eps_k(t)=&{\rm Tr} \left(\partial_\alpha H_kU_{\eps,k}(t)\, P_k\, U_{\eps,k}(t)^*\right) - {\rm Tr} \left(\partial_\alpha H_k P_k\right) - \eps t  {\rm Tr} \left(\partial_\alpha H_k \partial_\beta P_k\right) \\ & + i \eps \Tr(\partial_\alpha H_k (e^{-i t L_k} - 1) L_{k}^{+}  \partial_\beta P_k).
\end{align*}
We now expand the first term in the right-hand side of this equation. We set $\mu := 1 - \min \sigma(H)$,
$$H(s):=H_{k+se_\beta}+\mu,\quad A=\partial_\alpha H_k, \quad
I_{0}=[0,\frac 12(\lambda_{N,k}+\lambda_{N+1,k})+\mu],\quad P(s)=\1_{I_{0}}(H(s)).$$
It holds
$$H(s)=h_0+sh_1+\frac{s^2 |e_\beta|^{2}}2$$
 with $h_0=H_k+\mu$ and $h_1=\partial_\beta H_k$. The operators $h_0$, $h_1$ and $A$ are self-adjoint on $L^2_{\rm per}$ and we have $h_0 \ge 1$ and $h_1h_0^{-1/2}$ and $Ah_0^{-1/2}$ bounded.
  For $s$ in a neighborhood of zero, $P(s)= P_{k+se_\beta}$, so that $\dot{P}(0)=\partial_\beta P_k$.  Let $(U^\eps(t,t'))_{t,t' \in \R}$ be the propagator associated with the family $(H(\eps t))_{t \in \R}$ and $U^\eps(t):=U^\eps(t,0)$.
We have $U^\eps(t,t')=e^{-i\mu(t-t')} U_{\eps,k}(t,t')$ and $U^\eps(t)=e^{-i\mu t} U_{\eps,k}(t)$.
With these notations, we have
$$
\rho^\eps_k(t)={\rm Tr} \left(A U^\eps(t)\, P(0) U^\eps(t)^*\right) - {\rm Tr} \left(AP(0) \right) - \eps t  {\rm Tr} \left(A  \dot{P}(0) \right) + i \eps \Tr(A (e^{-i t L_0} - 1) L_0^{+}  \dot{P}(0)),
$$
where $L_0=L_k$ and $L_0^+=L_k^+$, and we focus on expanding the operator $U^\eps(t)\, P(0) U^\eps(t)^*$ close to $\eps=0$.

\begin{lemma}\label{lem:Dyson+}
We have
\begin{equation}\label{eq:Dyson+}
U^\eps(t)\, P(0) U^\eps(t)^* = P(0) +  \eps \left(t \dot{P}(0) - i \left( e^{-itL_0}-1 \right) (L_0^+\dot{P}(0))\right) + \Pi^\eps_2(t)
\end{equation}
with $\rho^\eps_k(t) = {\rm Tr} \left(A \Pi_2^\eps(t) \right)$.
Moreover, for $0\le t\le T$ and $0\le\eps\le1$, we have the bound~\eqref{eq:boundrho}
\begin{align*}
  |\rho^\eps_k(t)| \le C_T \eps^2.
\end{align*}

\end{lemma}

Lemma~\ref{lem:Dyson+} closes the proof of Proposition~\ref{prop:Dyson}.
\end{proof}


\begin{proof}[Proof of Lemma~\ref{lem:Dyson+}]
We deduce from the Dyson expansion, on $D(h_0^2)$ and in particular
on $\operatorname{Ran}P(0)$, that
$$
U^\eps(t)=U^0(t)+V^\eps(t)+W^\eps(t),
$$
where $c=|e_\beta|^2$, {$U^0(t)=e^{-ith_0}$} and
\begin{align*}
V^\eps(t) &= - i\eps \int_0^t U^0(t-t')  t' h_1 U^0(t') \, dt', \\
W^\eps(t) &=\eps^2\left(-i\frac{ct^3}{6}U^0(t)
-\int_0^t\int_0^{t'}U^\eps(t,t')t'(h_1+\eps ct'/2)
U^0(t'-t'')t''(h_1+\eps ct''/2)U^0(t'')\,dt''\,dt'\right).
\end{align*}
The products with $P(0)$ below have bounded finite-rank extensions.
This induces $U^\eps(t)\, P(0) U^\eps(t)^* = P(0) + \Pi^\eps_1(t)+\Pi^\eps_2(t)$
where
\begin{align*}
\Pi^\eps_1(t)&=V^\eps(t)P(0)U^0(t)^* + {\rm h.c.} = - i\eps \int_0^t t' U^0(t-t')  [h_1,P(0)] U^0(t-t')^* \, dt', \\
\Pi^\eps_2(t)&=V^\eps(t)P(0)V^\eps(t)^* + \left( W^\eps(t) P(0) (U^0(t)+V^\eps(t))^* + {\rm h.c.} \right) + W^\eps(t)P(0)W^\eps(t)^*.
\end{align*}

\medskip

We first analyze $\Pi^\eps_1(t)$ by computing
\begin{align*}
  U^0(t-t') [h_1,P(0)] U^0(t-t')^* &= - e^{-i(t-t') L_{0}} L_{0} \dot P(0)\\
  &= i\frac{d}{dt'} e^{-i(t-t') L_{0}} \dot P(0)\\
  &= \frac{d^{2}}{dt'^{2}} e^{-i(t-t') L_{0}} L_{0}^{+}\dot P(0),
\end{align*}
where we have used $\dot P(0) = L_{0}^{+}[P(0),h_{1}]$ and
$\dot P(0) = L_{0}^{+} L_{0} \dot P(0)$. Using integration by parts, we obtain
\begin{align*}
  \Pi^{\eps}_{1}(t) = \eps \left(t \dot{P}(0) - i \left( e^{-itL_0}-1 \right) (L_0^+\dot{P}(0))\right)
\end{align*}
and \eqref{eq:Dyson+} follows.


We now work on the bound~\eqref{eq:boundrho}. Since the limit
$\eps\to0$ is taken at fixed time, we only need estimates on $[0,T]$.
The graph and form norm equivalences give, uniformly in $k$,
\[
\|Ah_0^{-1/2}\|+\|h_1h_0^{-1/2}\|
+\|h_0^{1/2}h_1h_0^{-1}\|\le C.
\]
We also have
\[
h_0h_1h_0^{-2}=h_1h_0^{-1}
+i(e_\beta\cdot\nabla V)h_0^{-2}.
\]
The identity holds on $L^2_{\rm per}$: elliptic regularity gives
$h_0^{-2}L^2_{\rm per}\subset H^3_{\rm per}$, so $h_1h_0^{-2}$ takes
values in the domain of $h_0$. The last term is bounded because $\nabla V\in L^2_{\rm per}$ and
$h_0^{-2}$ maps $L^2_{\rm per}$ into $H^2_{\rm per}\subset L^\infty_{\rm per}$.
Thus $\|h_0h_1h_0^{-2}\|\le C$ as well.

Applying these bounds to the integrals defining $V^\eps$ and $W^\eps$,
and using~\eqref{eq:fixed_time_propagation}, we obtain
\begin{equation}\label{eq:Dyson_fixed_time}
\|h_0^{1/2}V^\eps(t)h_0^{-1}\|\le C_T\eps,
\qquad
\|h_0^{1/2}W^\eps(t)h_0^{-2}\|\le C_T\eps^2.
\end{equation}
For example, the term with two factors $h_1$ in $W^\eps$ is controlled by
\[
\|h_0^{1/2}U^\eps(t,t')h_0^{-1/2}\|
\,\|h_0^{1/2}h_1h_0^{-1}\|
\,\|h_0h_1h_0^{-2}\|,
\]
since $U^0$ commutes with $h_0$. The scalar terms are easier.
Using $P(0)^2=P(0)$, $\operatorname{Rank}P(0)=N$ and
$\|h_0^mP(0)\|\le(\mu_{\rm F}+\mu)^m$, we deduce
\[
\|AV^\eps(t)P(0)\|+\|V^\eps(t)P(0)\|\le C_T\eps,
\qquad
\|AW^\eps(t)P(0)\|+\|W^\eps(t)P(0)\|\le C_T\eps^2.
\]
Each trace in $\operatorname{Tr}(A\Pi_2^\eps(t))$ is therefore bounded
by $C_T\eps^2$. This proves~\eqref{eq:boundrho}, without using a lower
bound on the spectral gap.
\end{proof}


\section{Proof of the conductivity formula}\label{sec:semimet}

We prove here Theorem \ref{theo:semi-metal}. We first take the limit
$\eps\to0^+$ in Proposition~\ref{prop:Dyson}. The first-order Dyson
term and the bounds~\eqref{eq:Dyson_fixed_time} show that the derivative
of the fiber current is bounded uniformly in $k$ on each fixed time
interval. The remainder~\eqref{eq:boundrho} is uniform as well.
Together with Proposition~\ref{prop:null_current}, this gives
\begin{equation}\label{eq:linear_response_current}
j^{\rm LR}(t)=\frac1{4\pi^2}\int_{\mathcal B}j_k^{\rm LR}(t)\,dk,
\end{equation}
where, for $k\notin\mathcal S$,
\begin{equation}\label{eq:linear_response_fiber}
j_k^{\rm LR}(t)=t\partial_\alpha\partial_\beta\Tr(H_kP_k)
-i\Tr\left(\partial_\alpha H_k(e^{-itL_k}-1)L_k^+\partial_\beta P_k\right).
\end{equation}
The limit can also be passed under any integral over a fixed bounded
time interval. From now on we decompose the integral defining
$j^{\rm LR}(t)$ into several parts depending whether one integrates
far from the Dirac points or not.

\medskip

We introduce a small parameter $\delta>0$ controlling the size of the
neighborhood of the Dirac points, which is independent of $t$.
We decompose $\mathcal B$ as the disjoint union
\begin{align*}
  \mathcal B=\mathcal B_{\rm out}^{\delta}\; \cup  \, (\cup_{i=1}^M \mathcal B_{i}^{\delta})
\end{align*}
with
\[ \mathcal B_i^\delta= B(k_{i}, \delta),\]
where $\delta > 0$ is small enough so that
\[\mathcal B_{\rm out}^{\delta}\subset \{k\in\mathcal B,\;
  \lambda_{N,k} \leq \mu_F-c\delta\}\]
for some constant $c>0$. All these sets are understood modulo the
reciprocal lattice. Note that this decomposition is time-reversal
symmetric in the sense that
\[
  -\mathcal B_{\rm out}^{\delta} = \mathcal
B_{\rm out}^{\delta}\;\; \mbox{and}\;\;
  - (\cup_{i=1}^M \mathcal B_{i}^{\delta}) = (\cup_{i=1}^M \mathcal B_{i}^{\delta}).
\]
We let $t\to\infty$ with $\delta$ fixed, and finally let $\delta\to0$.
Write $\mathcal B_{\rm in}^\delta=\cup_{i=1}^M\mathcal B_i^\delta$.
On $\partial\mathcal B_i^\delta$, the vector surface element $ds$ points
into the ball (out of $\mathcal B_{\rm out}^\delta$).

\medskip

In the following analysis, we first treat the regions
$\mathcal B_{\rm out}^{\delta}$, where the linear response formula
has a non-zero gap larger than a constant times $\delta$. In the sets
$\mathcal B_{i}^{\delta}$, where the gap closes, we study the structure of the Taylor expansion of the Hamiltonian~$H_k$ close to the Dirac points and construct two-band reduced Hamiltonians~$H_{i,k}^{R}$. We reduce successively from the
Hamiltonian $H$ to the reduced Hamiltonian
$H_{i,k}^{R}$,
and finally to the
Dirac Hamiltonian
\[H_k^{D} = v\begin{pmatrix} 0 & k_1-ik_2 \\ k_1+ik_2 & 0\end{pmatrix},\qquad v>0,\]
for which we can explicitly
compute the current. Adding the contributions, we will obtain
\begin{align*}
\frac1t\int_0^t j^{\rm LR}(t')\,dt'
=\frac{M}{16}e_\alpha\cdot e_\beta+O(\delta)+O_\delta(t^{-1}).
\end{align*}
Finally, we will pass to the limits $t\to\infty$ and $\delta\to0$.

\subsection{Far from the Dirac points}
We set
\[
j^{{\rm LR},{\rm out}}(t):=
\frac1{4\pi^2}\int_{\mathcal B_{\rm out}^\delta}j_k^{\rm LR}(t)\,dk.
\]
We treat the two terms in~\eqref{eq:linear_response_fiber} separately.
\begin{itemize}
\item Using Stokes formula for the first term, we get
\begin{equation}\label{secondslot}
\int_{\mathcal B_{\rm out}^\delta}
t\partial_\alpha\partial_\beta\Tr(H_kP_k)\,dk
=t\sum_{n\le N}\int_{\partial\mathcal B_{\rm in}^\delta}
\partial_\alpha\lambda_{n,k}(ds\cdot e_\beta).
\end{equation}
Here derivatives of the sum of occupied eigenvalues are understood as
derivatives of $\Tr(H_kP_k)$, including at crossings within the occupied bands.
\item For the second term, the constant part vanishes when integrated
on the time-reversal symmetric set $\mathcal B_{\rm out}^\delta$.
Indeed, complex conjugation relates $k$ and $-k$; the trace
$\Tr(\partial_\alpha H_kL_k^+\partial_\beta P_k)$
is purely imaginary and odd. For the oscillatory part we use
\[
\frac1t\int_0^t e^{-i\omega t'}\,dt'
=\frac{e^{-i\omega t}-1}{-i\omega t}.
\]
In the eigenbasis representation~\eqref{eq:explicit_inverse_liouvillian},
only pairs $n\le N<m$ occur. On
$\mathcal B_{\rm out}^\delta$ their energy differences are bounded
away from zero for each fixed $\delta$. The finite number of occupied
states, the uniform bounds on $\partial_\alpha H_ku_{n,k}$ and
$\partial_\beta H_ku_{n,k}$, and Cauchy--Schwarz give absolute
summability of the products of matrix elements. Thus the time average
of the oscillatory part is $O_\delta(t^{-1})$.
\end{itemize}
We are left with
\begin{equation}\label{Bout}
\frac1t\int_0^t j^{{\rm LR},{\rm out}}(t')\,dt'
=\frac{t}{8\pi^2}\sum_{n\le N}
\int_{\partial\mathcal B_{\rm in}^\delta}
\partial_\alpha\lambda_{n,k}(ds\cdot e_\beta)
+O_\delta(t^{-1}).
\end{equation}
Here $O_\delta$ permits a constant depending on the fixed $\delta$.

\subsection{Close to Dirac points: reduction to the 2-band case}
\label{sec:dirac_reduc}

We set
\begin{equation}\label{eq:jepsin}
j^{{\rm LR},{\rm in}}(t):=\frac1{4\pi^2}
\sum_{i=1}^M\int_{\mathcal B_i^\delta}j_k^{\rm LR}(t)\,dk.
\end{equation}
The analytic two-band reduction described below, together with
Assumption~\ref{ass:semimetal}, shows that
\[
\partial_\alpha\partial_\beta\Tr(H_kP_k)=O(|k-k_i|^{-1})
\]
near $k_i$. This is integrable in dimension two. Since $j_k^{\rm LR}(t)$
is bounded at fixed $t$, the second term
in~\eqref{eq:linear_response_fiber} is also integrable on
$\mathcal B_i^\delta$.

\medskip

We treat the two terms in sequence.
\begin{itemize}
\item For the first, using Stokes formula, we get
\begin{align*}
\int_{\mathcal B_i^\delta}t\partial_\alpha\partial_\beta
\Tr(H_kP_k)\,dk
=-t\sum_{n\le N}\int_{\partial\mathcal B_i^\delta}
\partial_\alpha\lambda_{n,k}(ds\cdot e_\beta),
\end{align*}
so that the corresponding time-averaged term cancels the contribution
in~\eqref{Bout} from $\mathcal B_{\rm out}^\delta$.
To justify Stokes formula, first remove a ball of radius $r$ about $k_i$.
The extra boundary integral is $O(r)$, since
$\partial_\alpha\Tr(H_kP_k)=\Tr(\partial_\alpha H_kP_k)$ is bounded.
It therefore vanishes as $r\to0$.
\item For the second term, we use
\begin{align*}
&\Tr(\partial_\alpha H_k(e^{-itL_k}-1)L_k^+\partial_\beta P_k)
=\Tr(\partial_\alpha H_k(e^{-itL_k}-1)(L_k^+)^2[P_k,\partial_\beta H_k])\\
&=\sum_{n\le N}\sum_{m>N}
(e^{-it(\lambda_{n,k}-\lambda_{m,k})}-1)
\frac{\langle u_{n,k},\partial_\beta H_ku_{m,k}\rangle
\langle u_{m,k},\partial_\alpha H_ku_{n,k}\rangle}
{(\lambda_{m,k}-\lambda_{n,k})^2}-{\rm c.c.},
\end{align*}
with absolute convergence by the same Cauchy--Schwarz argument as above.
When $n\ne N$ or $m\ne N+1$, the denominators are bounded from below
independently of $\delta$. The constant term vanishes when integrated
over the time-reversal symmetric set $\cup_{i=1}^M\mathcal B_i^\delta$,
and the oscillatory term has time average $O(t^{-1})$.
\end{itemize}

Putting all the results of the previous two sections together, we get
\[
\frac1t\int_0^t j^{\rm LR}(t')\,dt'
=-\frac{i}{4\pi^2t}\sum_{i=1}^M\int_0^t I_i^R(\delta,t')\,dt'
+O_\delta(t^{-1}),
\]
with
\[
I_i^R(\delta,t)=\int_{\mathcal B_i^\delta}
(e^{-it(\lambda_{N,k}-\lambda_{N+1,k})}-1)
\frac{\langle u_{N,k},\partial_\beta H_ku_{N+1,k}\rangle
\langle u_{N+1,k},\partial_\alpha H_ku_{N,k}\rangle}
{(\lambda_{N+1,k}-\lambda_{N,k})^2}\,dk-{\rm c.c.}.
\]
At this stage of the proof, only two modes are involved in the formula
giving the current, namely the two modes that cross at the Fermi level.
Everything happens as for a two-band model that we now study.

\subsection{Close to the Dirac points: the local model}
We now are interested in the computation of $I_i^R(\delta,t)$.
In the following, we drop the index $i$ and assume without loss of
generality that $k_i=0$. We write $v=v_i$.

Hypothesis \eqref{hyp:semimetals_1} implies that, for $k$ small
enough, the Bloch Hamiltonian $H_k$ has exactly two eigenvalues
(counting multiplicities) close to $\mu_{\rm F}$. Consider an
arbitrary orthonormal basis $(v_0,w_0)$ of
$\Ran(P_{N+1,0} - P_{N-1,0})$. For all $k$ small
enough, we can construct an orthonormal basis $(v_k,w_k)$ of
$\Ran(P_{N+1,k} - P_{N-1,k})$ by L\"owdin
orthonormalization of
$((P_{N+1,k} - P_{N-1,k})v_0,(P_{N+1,k} -
P_{N-1,k})w_0)$, and set
$$
H_k^{R} = [v_{k} | w_{k}]^{*}
H_k[v_{k} | w_{k}] = \left( \begin{array}{cc} \langle v_k, H_k v_k \rangle &  \langle v_k, H_k w_k \rangle \\
\langle w_k, H_k v_k \rangle &  \langle w_k, H_k w_k \rangle \end{array} \right).
$$
It follows that the reduced Hamiltonian
\begin{align*}
  H_k\Big|_{\Ran(P_{N+1,k} - P_{N-1,k}) }
\end{align*}
is equivalent through a unitary transform that depends analytically on $k$
to the reduced $2 \times 2$ Hamiltonian
\begin{align}
  \label{eq:Hmatk}
  H_k^{R} = \sum_{p=0}^3 b^{p}(k) \sigma_{p},
\end{align}
where $$\sigma_{0} = {\rm Id}_{\C^{2}}, \;\; \sigma_{1}=\left( \begin{array}{cc} 0 & 1 \\ 1 & 0 \end{array} \right), \;\;\sigma_{2}=\left( \begin{array}{cc} 0 & -i \\ i & 0 \end{array} \right), \;\;\sigma_{3}=\left( \begin{array}{cc} 1 & 0 \\ 0 & -1 \end{array} \right)$$
are the Pauli matrices, and $(b^{p})_{p=0,1,2,3}$ are real-valued
analytic functions of $k$ in a neighborhood of $0$.

The matrix $H_k^{R}$ has eigenvalues
$$  \lambda_{\pm}(k) = b^{0}(k) \pm \sqrt{\sum_{p=1}^{3}b^{p}(k)^{2}}.
$$
It follows that
\begin{align*}
  b^{0}(k) &= \mu_{F} + O(|k|^{2}), \quad
  b^{p}(k) = v\langle q^{p}, k\rangle + O(|k|^{2}), \quad p=1,2,3,
\end{align*}
where the $(q^{p})_{p=1,2,3}$ are the rows of a $3 \times 2$
matrix $Q$ with orthonormal columns, so that
\begin{align}\label{def:HR}
  H_k^{R} = \mu_{F} + v (Q k) \cdot \sigma + O(|k|^{2}).
\end{align}
Let $R \in {\rm SO}(3)$ be a rotation matrix that maps
${\rm Ran}(Q)$ to ${\rm Span}(e^0_{1},e^0_{2})$, where $(e^0_{1},e^0_{2},e^0_{3})$ is the canonical basis of $\R^3$. Write $R$ as a rotation through an angle $\theta$ about a unit vector $n$.
The unitary matrix
\[
U=\cos(\theta/2)\mathrm{Id}_2-i\sin(\theta/2)n\cdot\sigma
\]
satisfies, by the Pauli matrix identities,
\begin{align*}
  R_{pq} = \frac 1 2 \Tr(\sigma_{p} U \sigma_{q} U^{*}).
\end{align*}
It follows that
\begin{align*}
  \Tr(\sigma_{3} U H_k^{R} U^{*}) &= v\sum_{q=1}^{3} (Q k)_{q} \Tr(\sigma_{3} U \sigma_{q} U^{*}) + O(|k|^{2})\\
  &= 2 v\langle e^0_3,RQk\rangle+O(|k|^2)\\
  &= O(|k|^{2}).
\end{align*}
Up to a unitary transform, we can therefore assume $Q$ to be
a $2\times 2$ matrix in~\eqref{def:HR}.

\subsection{The two-band case: reduction to the Dirac Hamiltonian}
\paragraph{Reduction to $H_k^{R}$}
For $k \neq 0$, let $\lambda^{R}_{\pm,k}$ be the larger and smaller
eigenvalues of $H_k^{R}$ respectively, and $u^{R}_{\pm,k}$
associated orthonormal eigenvectors in $\C^2$. We have
$\lambda_{-,k}^{R} = \lambda_{N,k}, \lambda_{+,k}^{R} =
\lambda_{N+1,k}$, and
\begin{align*}
  [v_{k} | w_{k}]
u_{+,k}^{R} = e^{i\theta_{+}(k)}u_{N + 1,k},\qquad[v_{k} | w_{k}]
u_{-,k}^{R} = e^{i\theta_{-}(k)}u_{N,k}
\end{align*}
for some phases $\theta_{\pm}(k) \in \R$. We have
\begin{align*}
  \partial_\alpha H_k^{R}
  &= [v_{k} | w_{k}]^{*}
\partial_\alpha H_k[v_{k} | w_{k}]+ \partial_\alpha [v_{k} | w_{k}]^{*} [v_{k} | w_{k}] H_k^{R} + H_k^{R}  [v_{k} | w_{k}]^{*}\partial_\alpha [v_{k} | w_{k}]
\\
&= [v_{k} | w_{k}]^{*}
\partial_\alpha H_k[v_{k} | w_{k}] + O(|k|)
\end{align*}
where we have used for the first line that $H_k$ commutes with
$P_{N+1,k} - P_{N-1,k} = [v_{k} | w_{k}]
[v_{k} | w_{k}]^{*}$, and for the second that $H_k^{R} = \mu_{F} {\rm Id}_{2} + O(|k|)$ and
$\partial_\alpha \Big([v_{k} | w_{k}]^{*} [v_{k} | w_{k}] \Big) =
\partial_\alpha {\rm Id}_{2} = 0$. We therefore obtain
\begin{align*}
  \langle u^{R}_{+,k},\partial_\alpha H_k^{R}u^{R}_{-,k}\rangle&= e^{-i(\theta_{+}(k)-\theta_{-}(k))}\langle u_{N+1,k},\partial_\alpha H_ku_{N,k}\rangle + O(|k|)\\
  \langle u^{R}_{-,k}, \partial_\beta H_k^{R} u^{R}_{+,k} \rangle&= e^{+i(\theta_{+}(k)-\theta_{-}(k))}\langle u_{N,k},\partial_\beta H_ku_{N+1,k}\rangle + O(|k|).
\end{align*}
Since $(\lambda^{R}_{+,k} - \lambda^{R}_{-,k})$ is bounded from below
by a constant multiple of $|k|$, it follows that
\begin{equation}
  \label{eq:IRdt}
  I^{R}(\delta,t) = \int_{B(0,\delta)} (e^{-i t (\lambda^{R}_{-,k} - \lambda^{R}_{+,k})} - 1)\frac{\langle u^{R}_{-,k}, \partial_\beta H_k^{R} u^{R}_{+,k} \rangle\langle u^{R}_{+,k},\partial_\alpha H_k^{R}u^{R}_{-,k}\rangle}{(\lambda^{R}_{+,k} - \lambda^{R}_{-,k})^{2}}dk - {\rm c.c.} + O(\delta).
\end{equation}

\paragraph{Reduction to $H_k^{D}$}
By standard results of perturbation theory \cite[Chapter XII]{ReedSimon4} applied to
$H_k^{R} = \mu_{F} + v (Q k) \cdot \sigma + O(|k|^{2})
$ with gap greater than a constant multiple of $|k|$,
\[
  \lambda^{R}_{\pm,k} = \mu_F+\lambda^{Q}_{\pm,k} +
  O(|k|^{2})\;\;\mbox{and}\;\;   u^{R}_{\pm,k} = u^{Q}_{\pm,k} + O(|k|)
\]
where the superscript $Q$ refers to eigenvalues and appropriately
chosen orthonormal eigenvectors of the Hamiltonian
\begin{align*}
  H_k^{Q} = v(Qk) \cdot \sigma.
\end{align*}
We claim that, uniformly in $t\ge0$,
\begin{align*}
  I^{R} (\delta,t) &= I^{Q}(\delta,t) + O(\delta),
\end{align*}
where $I^{Q}(\delta,t)$ is defined similarly to
\eqref{eq:IRdt} as
\begin{align*}
  I^{Q}(\delta,t) = \int_{B(0,\delta)}  (e^{-i t (\lambda^{Q}_{-,k} - \lambda^{Q}_{+,k})} - 1)\frac{\langle u^{Q}_{-,k}, \partial_\beta H_k^{Q} u^{Q}_{+,k} \rangle\langle u^{Q}_{+,k},\partial_\alpha H_k^{Q}u^{Q}_{-,k}\rangle}{(\lambda^{Q}_{+,k} - \lambda^{Q}_{-,k})^{2}}dk - {\rm c.c.}.
\end{align*}


To prove the uniformity in $t$, write $k=r\omega(\theta)$ and set
\[
\Delta(r,\theta)=\lambda^R_{+,k}-\lambda^R_{-,k},\qquad
z(r,\theta)=\langle u^R_{-,k},\partial_\beta H_k^Ru^R_{+,k}\rangle
\langle u^R_{+,k},\partial_\alpha H_k^Ru^R_{-,k}\rangle.
\]
Analyticity of the reduced Hamiltonian gives
$\Delta=2vr+O(r^2)$ and $\partial_r\Delta=2v+O(r)$,
uniformly in $\theta$. For $H_k^Q$, the velocity product is a real
function $a(\theta)$ independent of $r$: the Pauli vectors of
$H_k^Q$, $\partial_\alpha H_k^Q$ and $\partial_\beta H_k^Q$ lie in the
same plane. The eigenvector estimates above therefore give
$\operatorname{Re}z=a(\theta)+O(r)$ and
$\operatorname{Im}z=O(r)$.
After subtracting the complex conjugate in~\eqref{eq:IRdt}, the radial
integrand is
\[
2i\bigl(\sin(t\Delta)\operatorname{Re}z
+(\cos(t\Delta)-1)\operatorname{Im}z\bigr)\frac{r\,dr}{\Delta^2}.
\]
The imaginary part of $z$ contributes $O(\delta)$, uniformly in $t$.
For the remaining term, perform the change of variable
$\rho=\Delta(r,\theta)/(2v)$. We have
\[
\frac{(\operatorname{Re}z)r\,dr}{\Delta^2}
=\left(\frac{a(\theta)}{4v^2\rho}+O(1)\right)d\rho,
\qquad \rho(\delta,\theta)=\delta+O(\delta^2).
\]
The $O(1)$ term integrates to $O(\delta)$, since the sine is bounded.
Changing the upper endpoint back to $\delta$ also costs $O(\delta)$:
the interval has length $O(\delta^2)$ and $\rho$ is comparable to
$\delta$ there. The remaining integral is exactly $I^Q(\delta,t)$.
This proves the claim without expanding the oscillatory phase in $t$.

We perform the change of variable $k' = Qk$ (recall that $Q$ is orthogonal) and obtain
\begin{align*}
I^{Q}(\delta,t) = (Qe_\beta)^T I^{D}(\delta,t)(Qe_\alpha),
\end{align*}
where the coefficients $I_{ij}^{D}(\delta,t) $ of the $2\times 2$ matrix  $I^{D}(\delta,t) $ are given by
\begin{align*}
  I_{ij}^{D}(\delta,t) = \int_{B(0,\delta)} (e^{-i t (\lambda^{D}_{-,k} - \lambda^{D}_{+,k})} - 1)\frac{\langle u^{D}_{-,k}, \partial_{k_j} H_k^{D} u^{D}_{+,k} \rangle\langle u^{D}_{+,k},\partial_{k_i} H_k^{D}u^{D}_{-,k}\rangle}{(\lambda^{D}_{+,k} - \lambda^{D}_{-,k})^{2}}dk - {\rm c.c.}
\end{align*}
and the superscript $D$ refers to the Dirac Hamiltonian
\begin{align*}
  H_k^{D} = v k \cdot \sigma.
\end{align*}


\paragraph{The Dirac Hamiltonian $H_k^{D}$}
We finish by computing $I^{D}(\delta,t)$ explicitly. Let
$k=r(\cos \theta, \sin \theta)$. We have
\[  \lambda^{D}_{\pm,k} = \pm v r,\;\;
  u^{D}_{+,k} = \frac 1 {\sqrt 2}
  \begin{pmatrix}
    1\\e^{i\theta}
  \end{pmatrix}\;\;\mbox{and}\;\;
  u^{D}_{-,k}= \frac 1 {\sqrt 2}
  \begin{pmatrix}
    -e^{-i\theta}\\1
  \end{pmatrix}.
  \]
By an explicit calculation, we obtain
\begin{align*}
  \int_{0}^{2\pi}(e^{-i t (\lambda^{D}_{-,k} - \lambda^{D}_{+,k})}-1)\frac{\langle u^{D}_{-,k}, \partial_{k_i} H_k^{D} u^{D}_{+,k} \rangle\langle u^{D}_{+,k},\partial_{k_j} H_k^{D} u^{D}_{-,k}\rangle}{(\lambda^{D}_{+,k} - \lambda^{D}_{-,k})^{2}} \, d\theta - \text{c.c.} &= i\pi \frac 1 {2r^{2}}\sin(2vrt) \delta_{ij}.
\end{align*}
It follows that
\begin{align*}
  \frac 1 t \int_{0}^{t}I_{ij}^{D}(\delta,t') \, dt' &= \frac {i\pi}{2}\delta_{ij} \frac1t\int_0^t \int_{0}^{\delta}\frac{\sin(2vrt')}{r} \, dr \, dt' = \frac{i\pi}{4 t}\delta_{ij}\int_{0}^{\delta}\frac{1 -  \cos(2vrt)}{vr^{2}} \, dr\\
  &=\frac{i\pi}{4} \delta_{ij} \int_{0}^{\delta vt}\frac{1-\cos(2r')}{(r')^2} \, dr' = \frac{i \pi^{2}}{4}\delta_{ij} + O\left((\delta t)^{-1}\right).
\end{align*}
We finally get by summing all the estimates
\begin{align*}
  \frac 1{t} \int_{0}^{t} j^{\rm LR}(t') \, dt' = \frac {M} {16} e_\alpha \cdot e_\beta   +O(\delta)+O_\delta(t^{-1}).
\end{align*}
Since $Q$ is orthogonal, first letting $t\to\infty$ and then
$\delta\to0$ proves the result.


\appendix


\section{Proofs of two technical lemmata}

\subsection{Proof of Lemma~\ref{FarisLavine}}\label{sec:FL_lemma}
\begin{proof}
We use the proof of the Faris--Lavine theorem in \cite{ReedSimon2}. It consists in verifying the following two hypotheses of \cite[Theorem X.37]{ReedSimon2}.  Let $A = -\frac12\Delta + W+V$ and $K=A+2c|x|^2 +b$, where $b\in\R$ will be specified below:
\begin{equation}
\label{HypX37.1}\textit{there exists $h$, such that for any } \phi \in \mathcal{C},
\quad \|A\phi\|\leq h \|K\phi\|;
\end{equation}
\begin{equation}
\label{HypX37.2}
\textit{for some $\ell$, for any }  \phi \in \mathcal{C},\quad |(A\phi,K\phi) - (K\phi,A\phi)| \leq \ell \|K^\frac{1}{2} \phi\|^2.
\end{equation}

By hypothesis \textit{3} in Lemma~\ref{FarisLavine} and the conditions on $W$, it is possible to choose $b$ so that $K\geq 1$.
As quadratic forms on $\mathcal{C}$,
\[K^2 = (A+b)^2 + 4c \sum_{j=1}^2 x_j (A+b+c|x|^2)x_j - 4c.  \]
Hypotheses \textit{1} and \textit{3} guarantee that $A+b+c|x|^2$ is bounded below.
Hence, increasing the value of~$b$ if necessary to make this operator positive, we have
\[ \|(A+b)\phi\|^2_{L^2} \leq \|K\phi \|^2_{L^2} +8c\|\phi \|^2_{L^2},\]
which proves (\ref{HypX37.1}).

\medskip

For (\ref{HypX37.2}), we observe that
\[ \pm i [A,K] = \pm 2c(x \cdot (-i\nabla) + (-i\nabla)\cdot x) \leq 2 c \left((-\Delta) +|x|^2\right) \leq \ell K, \]
where we have used
\[(-\Delta) + |x|^2 \pm (x\cdot(-i\nabla) + (-i\nabla)\cdot x) = (-i\nabla  \pm x)^2 \geq 0 \]
and
\[ K = \left(\frac a 2 (-\Delta) + V\right) + (W+c|x|^2) + \frac{1-a}{2}(-\Delta) + c|x|^2 +b \geq e((-\Delta)+|x|^2), \]
where $e = \min(c,\frac{1-a}{2})>0$ and where $b$ is chosen so that
$$b- f +\min \,\sigma\left(\frac a 2 (-\Delta) + V\right) \geq 0.$$
This proves  (\ref{HypX37.2}).
Hence $A$ is essentially self-adjoint on $\mathcal{C}$.
\end{proof}

\subsection{Proof of Lemma~\ref{lemmaUnitaryPropagator}}\label{sec:unitaryprop}

\begin{proof}By the Kato-Rellich theorem, for any $0\leq t\leq T$,
  $H(t)$ is self-adjoint on $L^2_{\rm per}$ with domain~$H_{\rm per}^{2}$, and
  bounded below. We will show that there exists $\mu>0$ so that the
  graph norm of $(H(t)+\mu)$ for any $0\leq t \leq T$ is equivalent to
  the $H_{\rm per}^{2}$-norm. Together with the assumed $C^1$ dependence
  on $t$, this verifies the hypotheses of the propagator theorem
  \cite[Theorem X.70]{ReedSimon2}, and proves
  Lemma~\ref{lemmaUnitaryPropagator}.

\medskip

We have for any $\mu>0$, $0\leq t\leq T$ and $\phi \in H_{\rm per}^{2}$,
\[ \|(H(t)+\mu)\phi \|_{L^{2}_{\rm per}} \leq (1+a)\|H_{0} \phi \|_{L^{2}_{\rm per}} + (b+\mu) \|\phi\|_{L^{2}_{\rm per}} \leq (1+a+b+\mu) \|\phi\|_{H_{\rm per}^{2}},\]
and so the graph norm is controlled by the $H_{\rm per}^{2}$-norm.

\medskip

For the other inequality, we relate the resolvent of $H(t)$ to that of $H_{0}$ by a bounded operator, with bounded inverse.
Notice that, for any $\mu>0$, since $H_{0}$ is positive,
\[\forall \; 0\leq t\leq T, \quad (H(t)+\mu) = (1+H_{1}(t)(H_{0}+\mu)^{-1})(H_{0}+\mu).\]
Furthermore,
\[\forall \; 0\leq t\leq T, \quad\|H_{1}(t)(H_{0}+\mu)^{-1}\| \leq
  a\|H_{0}(H_{0}+\mu)^{-1}\| +b\|(H_{0}+\mu)^{-1}\| \leq a
  +\frac{b}{\mu}.\]
and so, for $\mu > \frac{b}{1-a}$, the operator
$1+ H_{1}(t)(H_{0}+\mu)^{-1}$ is bounded and invertible with bounded
inverse in $L^{2}_{\rm per}$. Therefore $(H(t)+\mu)^{-1}$ is bounded
from $L^2_{\rm per}$ to $H_{\rm per}^{2}$, which means there exists $C>0$ such
that, for any $\phi \in H_{\rm per}^{2}$ and $0\leq t\leq T$,
\[  \|\phi\|_{H_{\rm per}^{2}} = \|(H(t)+\mu)^{-1}(H(t)+\mu)\phi\|_{H_{\rm per}^{2}} \leq C\|(H(t)+\mu)\phi\|_{L^{2}_{\rm per}}, \]
which concludes the proof.
\end{proof}


\bibliographystyle{plain}
\bibliography{refs}
\end{document}
